\documentclass[11pt]{article}
\usepackage[margin=1in]{geometry}
\usepackage{amsmath,amssymb,amsthm}
\usepackage[T1]{fontenc}
\usepackage[hidelinks]{hyperref}
\newcommand{\SwapReg}{\operatorname{SwapReg}}
\newcommand{\norm}[1]{\left\lVert#1\right\rVert}
\newcommand{\rn}[1]{\norm{#1}_{\infty,1}}
\newcommand{\Sh}{\mathsf S}
\newcommand{\Dif}{\mathsf D}
\newcommand{\Fil}{\mathsf F}
\newcommand{\Res}{\mathsf R}
\newcommand{\Pred}{\mathsf H}
\newtheorem{theorem}{Theorem}
\newtheorem{lemma}[theorem]{Lemma}
\newtheorem{proposition}[theorem]{Proposition}
\title{prb_2a3012883e274a1384504ef3c9221892: a finite-bit transfer for constant swap regret}
\author{OpenAI GPT 6 Astra Pro}
\date{18 October 2026}
\begin{document}
\maketitle

\section*{Problem}
\textbf{prb_2a3012883e274a1384504ef3c9221892 (live submission key \texttt{ECON-3}).}
Finite-bit uncoupled learning with constant swap regret.
\href{https://we-solve.github.io/problems/finite-bit-uncoupled-learning-with-constant-swap-regret/}{We-Solve problem page.}

\section*{Title}
Rational transition matrices, exact stationary distributions, and a
summable-error transfer for constant swap regret.

\section*{Claim type}
\textbf{Full solution (affirmative), using the cited analytic estimates
of Mai, version 2.}

\section*{Scope}
We give a deterministic finite-bit learner and prove that it answers
the exact catalogue question affirmatively, preserving its optional
dimension rate up to an absolute constant.  The proof controls the
learner's \emph{actual} self-play history, including accumulated
feedback and state-rounding errors.  It uses Mai's higher-order
prediction estimates as cited results, and proves that they apply
with the bounded initial-score perturbation needed by this
construction.  The proof supplies the numerical and memory bounds
excluded from the cited exact-real implementation.

\section*{Assumptions and limitations}
The game, oracle, feedback restriction, regret benchmark, and resource
requirements are exactly those in the attached statement.  The proof
counts oracle reply lengths and all subsequent bit operations, while
the oracle's internal expectation calculation remains excluded as in
the problem.  All outputs and stored states are rational.  No realized
sampled-action guarantee, fixed-precision guarantee, or adversarial
guarantee is asserted.  The payoff bit bound $L$ is not needed by the
construction beyond the problem's provision of the oracle.

The principal literature input is Mai's version 2, whose title is
\emph{Constant Swap Regret in General-Sum Games via Two-Scale Higher-Order
Optimism}.  Version 1, cited in the question, has a different title.
Both the model and the explicit finite-bit exclusion in Appendix H were
checked in version 2 \cite{mai}.  The cited analytic estimates and
their hypotheses are identified below; they are not independently
reconstructed in full.  This is a proposed mathematical proof using
those results, with separate-agent checks of the new argument.  It is
not a formally verified proof or an externally refereed acceptance.

\section*{Write-up}
\subsection*{1. Notation and a rational choice of public parameters}
Write $n=N$ for the number of original players.  Let
\[
 s_i=2^{\lceil\log_2\max\{2,m_i\}\rceil},\qquad s=\max_i s_i\leq2K.
\]
Replace each original action list by $s_i$ virtual actions, using
duplicates of existing actions to fill the list.  Fix a public
surjection $p_i:[s_i]\to[m_i]$.  The virtual game has payoffs
\[
 \widetilde u_i(a_1,\ldots,a_n)
 =u_i(p_1(a_1),\ldots,p_n(a_n)).
\]
If the virtual mixed strategy is $y_i$, the learner outputs
$x_{ia}=\sum_{b:p_i(b)=a}y_{ib}$.  The original payoff-vector reply,
copied onto the virtual duplicates, is a reply for this virtual game
at the same precision.  No extra information is requested.

For an original map $\varphi:[m_i]\to[m_i]$, choose for every original
action a fixed virtual representative.  Mapping each virtual action
$b$ to the representative of $\varphi(p_i(b))$ gives a virtual
remapping with exactly the original map's cumulative gain.  Hence
\begin{equation}\label{eq:padding}
 \SwapReg_i^{\mathrm{original}}(T)
 \leq\SwapReg_i^{\mathrm{virtual}}(T).
\end{equation}

For the rest of the proof, $r_i^t$ denotes the true virtual gain matrix,
\[
 r_{i,a,b}^t=y_{i,a}^t
 (v_{i,b}^t-v_{i,a}^t),\qquad
 \rn{w_i}=\sum_{a=1}^{s_i}\max_b|w_{i,a,b}|.
\]
Thus $\rn{r_i^t}\leq1$ and
\[
 \SwapReg_i(T)=\sum_a\max_b\sum_{t\leq T}r_{i,a,b}^t\geq0.
\]

Use Mai's public normalizer and predictor parameters, with all
unqualified logarithms in their base-two convention:
\[
 M=\sum_i s_i^2,\quad A_i=s_i\log_2s_i,\quad A=\sum_iA_i,
 \quad\ell=\ell_0+\lceil\log_2(2+n+M)\rceil,\quad k=\ell+1,
\]
\[
 \delta=\frac1{n(M+\ell+2)^2},\qquad
 J=\lceil\log_2(1/\delta)\rceil,\qquad
 \beta_i=\frac{\log_2s_i}{2^{10}s_iJ}.
\]
Let $d_i$ be the least power of two at least
$\lceil\log_2(1/\beta_i)\rceil$.  These parameters are rational or
integer: this is the reason for the virtual-action padding.
For later use, $M\geq s_i^2$ implies
$\delta^{-1}=n(M+\ell+2)^2\geq s_i^4$, hence
$J\geq4\log_2s_i$ and $s_i\beta_i\leq2^{-12}$.
The base-two convention also appears in the source's identity
$\ln m=(\ln2)\log m$ in its potential estimate.

Fix a sufficiently small positive \emph{rational} universal $c$ and
a sufficiently large universal integer $\ell_0$, as in the proof of
Lemma~\ref{lem:history} below.  Put $g=c/k^{5/2}$ for the analysis only.  Choose
$\eta_i=2^{-h_i}$, where $h_i$ is the smallest nonnegative integer with
\[
 2^{-2h_i}\leq\frac{c^2A_i}{k^5A}.
\]
This choice uses only integer comparisons after clearing denominators.
It satisfies
\begin{equation}\label{eq:rates}
 \frac12g\sqrt{A_i/A}<\eta_i\leq g\sqrt{A_i/A},
 \qquad\sum_i\eta_i^2\leq g^2.
\end{equation}
In particular $\eta_i^{-1}$ and all public parameter bit lengths are
polynomially bounded in $n,K$.  The universal constants can be fixed
once and hardcoded.  The question asks for the existence of fixed
polynomials and does not require optimized numerical values of these
universal constants.  The proof of Mai's Theorem 2.1 explicitly treats
arbitrary positive rates satisfying its equation (4.1),
$\sum_i\eta_i^2\leq g^2$.  It therefore applies to these dyadic rates;
the exact default square-root allocation is not required.

\subsection*{2. The exact analytic interface}
For a row vector $z\in\mathbb R^{s_i}$, define
\[
 u(v)=\frac{\sqrt{v^2+4}-v}{2},\quad
 t_+(v)=u(v)^{-1},\quad
 f_i(v)=\beta_i
 \frac{(1+t_+(v)/d_i)^{d_i}}
 {1+u(v)+(\delta/4)u(v)^2}.
\]
There is a unique shift $\mu$ with
$\sum_b f_i(z_b-\mu)=1$.  Write
$q_{i,b}(z)=f_i(z_b-\mu)$.  Apply this map separately to every row
of a score matrix $z_i$, producing a positive transition matrix
$Q_i(z_i)$.  Let $y_i(z_i)$ be its unique stationary distribution.
The map $\Gamma_i(z)$ is the true gain matrix obtained in the fixed
virtual game when all players use these stationary distributions.

Let $\Sh$ denote backward shift on sequences with zero prehistory,
$\Dif=I-\Sh$, and
\[
 \Res_r=(I-\rho_r\Sh)^{-1},\qquad
 \Fil_r=\Res_r\Dif,\qquad\rho_r=\frac r{r+1},\qquad
 \Pred=\Dif^3\Fil_\ell^\ell\Fil_M^M.
\]
The exact predictor is $(I-\Pred)r$.  It uses only past gains because
the current-input coefficient of $\Pred$ is one.

\begin{lemma}[Robust history estimate from Mai's analytic lemmas]\label{lem:history}
For the normalizer and public parameters above, there are universal
constants, independent of the game and horizon, with the following
property.  Suppose a deterministic infinite effective-score history
satisfies, for fixed matrices $b_i$ with $\rn{b_i}\leq1/16$,
\begin{align}
 r_i^t&=\Gamma_i(z^t),\label{eq:truehistory}\\
 z_i^t&=\eta_i\left(b_i+\sum_{s<t}r_i^s
                 +(I-\Pred)[r_i]^t+\varepsilon_i^t\right),
 \quad\rn{\varepsilon_i^t}\leq(t+1)^{-2}.
 \label{eq:history}
\end{align}
Then, for sufficiently small universal $c$ and sufficiently large
universal $\ell_0$, every finite horizon satisfies
\begin{equation}\label{eq:inputconclusion}
 \SwapReg_i(T)\leq \frac{4A_i}{\eta_i}.
\end{equation}
The fixed $b_i$ need not be known to a learner and need not be a causal
function of its observations.
\end{lemma}
\begin{proof}
Fix a finite horizon $T$.  We first specify the source's energies and
prove that its prediction estimate applies to the stated history;
we then account for the nonzero initial score and choose universal
parameters that close the estimates.  Put
\[
 \xi_i=\Pred[r_i]-\varepsilon_i,\qquad
 \Theta_i^t=\eta_i\left(b_i+\sum_{s\leq t}r_i^s\right).
\]
Then $z_i^t=\Theta_i^t-\eta_i\xi_i^t$, and, for $t\geq2$,
\begin{equation}\label{eq:scoreincrement}
 z_i^t-z_i^{t-1}
 =\eta_i(r_i^t-\xi_i^t+\xi_i^{t-1}).
\end{equation}
Mai's Lemma G.1 gives $\norm{\Pred}_{\mathrm{ker},1}<96$, so
$\rn{\xi_i^t}\leq97$.

\paragraph{Prediction estimate and its history hypotheses.}
For one normalized row at score $z$, define
\[
 H_b=f_i'(z_b-\mu),\qquad
 \sigma_b=\frac{H_b}{\sum_cH_c},\qquad
 \mathcal K_z(w)=\sum_bH_b
 \left(w_b-\sum_c\sigma_cw_c\right)^2.
\]
These are Mai's equation (4.2), with $H_b$ denoting the source's
curvature weight $h_b$ and $\mathcal K$ its row quadratic form.
Define
\[
 S_i=\frac{\eta_i^2}{8}\sum_{t=1}^T\sum_a
       \mathcal K_{z_{i,a}^t}(r_{i,a}^t),\qquad
 \omega_i=C_\omega\frac{s_i-1}{\beta_i},
\]
\[
 Y^2=\sum_i\omega_iS_i,\quad V=\sum_i\omega_iA_i,
 \quad\Lambda^2=\sum_i\omega_i\eta_i^2,\quad
 E^2=\max_i\sum_{t=1}^T\rn{\xi_i^t}^2.
\]
Here $\|w\|_{\infty,1;T}=(\sum_{t=1}^T\|w^t\|_{\infty,1}^2)^{1/2}$.
The constant $C_\omega\geq1$ is the fixed universal curvature-comparison
constant of the source's Appendix A.3.  These definitions agree
with Section 4 immediately after equation (4.2) and before Lemma 4.1.
Mai's Lemma 4.2, equation (4.4), and its proof in Sections 4.3--4.7
and Appendices C--G give
\begin{equation}\label{eq:prediction}
 E\leq C_0+C_1gE+\tau_\ell Y,
 \qquad\tau_\ell=C_Pk^2(C_Egk^{5/2})^\ell.
\end{equation}

The source states its implementation with zero initial scores.  We
verify that this restriction does not enter the cited prediction
calculation.  The proof uses \eqref{eq:scoreincrement} only when
$t\geq2$.  Section 4.4 assigns $t=1$ the boundary impulse $r_i^1$,
whose row norm is at most one independently of $z^1$.  Its later early
expressions are gated by $t>p$, and its late expressions by
$t>\ell+j$.  More explicitly, differentiating a depth-$p$ early
expression occurs only at $t\geq p+2$.  Its lag indices are at most
$p$, so every increment index $t-a$ is at least two.  The late
calculation has the identical property with $p=\ell+j$.  The remaining
first-time term is a bounded boundary observable.  Thus no derivative
identity involving an initial, nonpositive-time score is needed.
The derivative estimates of Lemma C.1 and tree-score estimates of
Lemma D.1 are uniform over all real score vectors.  The late-window
comparison in Lemma F.1 uses differences of $\Theta$, in which the
fixed $b_i$ cancels.  In particular its hull diameter remains at most
$\eta_i(\ell+j+195)$, exactly as in that lemma.

The remaining entry of the perturbation into the prediction proof is
\[
 E\leq1+\max_i\norm{\Pred r_i}_{\infty,1;T},
\]
which follows from
$\sum_{t=1}^T\rn{\varepsilon_i^t}^2\leq
\sum_{t\geq1}(t+1)^{-4}<1$.
Its sign and computational origin do not occur in the estimate.
Thus the same early and late one-difference closures, Lemmas 4.3 and
4.4, apply with the same constants.  The same causal filter bounds
in Lemma G.1 and the finite assembly in Appendix G.2 prove
\eqref{eq:prediction} for the present history.  This establishes the
required extension of the cited analytic estimate; no history at
time zero has been added as an extra assumption.

\paragraph{Potential estimate with the initial score.}
Mai's row potential $\Phi_i$, defined in equation (4.7) and analyzed
in Lemma B.2, has gradient $q_i$, hence is
1-Lipschitz in the sup norm; it also satisfies
$\Phi_i(z)\geq\max_bz_b$ and $\Phi_i(0)\leq2\log_2s_i$.
The same one-round Taylor inequality as in the source has the same
two displacements $\eta_i\xi_i^t$ and
$\eta_i(\xi_i^t-r_i^t)$, including at $t=1$.  Stationarity gives
$\sum_{a,b}q_{i,a,b}^tr_{i,a,b}^t=0$ exactly.  The two endpoint
inequalities are now
\[
 \sum_a\Phi_i(\Theta_{i,a}^0)
 \leq2A_i+\eta_i\rn{b_i},\qquad
 \sum_a\Phi_i(\Theta_{i,a}^T)
 \geq\eta_i\SwapReg_i(T)-\eta_i\rn{b_i}.
\]
Consequently,
\begin{equation}\label{eq:biasedenergy}
 \eta_i\SwapReg_i(T)+S_i
 \leq2A_i+2\eta_i\rn{b_i}+\eta_i^2E^2.
\end{equation}
Because $A_i\geq2$, $\eta_i\leq1$, and $\rn{b_i}\leq1/16$,
this gives
\[
 Y^2\leq3V+\Lambda^2E^2,\qquad
\Lambda^2\leq g^2V.
\]
For the last inequality, use
$\Lambda^2\leq(\max_i\omega_i)\sum_i\eta_i^2\leq g^2V$.
This is the general-rate inequality of Mai's Lemma 4.1, rather than
the special identity available for the exact default allocation.

\paragraph{Universal parameters and absorption.}
Fix universal upper bounds for the analytic constants as in
Appendix A.3.  Choose a positive rational $c$ small enough that
\[
 C_Ec\leq\frac1{16},\qquad
 (C_1+1)c\leq\frac12,\qquad
 c^2[2(C_0+1)]^2\leq1,
\]
and small enough for the local potential comparisons in Section 4.2.
These finitely many bounds depend only on the cited universal
constants, not the game, dimensions, or horizon.

The source's dimension estimates also justify the enlarged $3V$
budget uniformly when $\ell_0$ changes.  Set
$v=\lceil\log_2(2+n+M)\rceil$, so $k=\ell_0+v+1$.
Appendix A.3 gives $J\leq5k$.  The definitions show directly that
\[
 V=2^{10}C_\omega J\sum_i s_i^2(s_i-1)
 \leq2^{13}C_\omega kM^{3/2}.
\]
Use $M\leq2^v$, $k\leq(\ell_0+1)(v+1)$, and
$C_Egk^{5/2}=C_Ec\leq1/16$ to obtain
\begin{align*}
 \tau_\ell\sqrt{3V}
 &\leq C_P\sqrt{3\cdot2^{13}C_\omega}
        k^{5/2}16^{-\ell}M^{3/4}\\
 &\leq C_P\sqrt{3\cdot2^{13}C_\omega}
 (\ell_0+1)^{5/2}16^{-\ell_0}
 (v+1)^{5/2}2^{-13v/4}.
\end{align*}
Here $v\geq3$ and $(v+1)^{5/2}2^{-13v/4}\leq1$.
Choose one universal $\ell_0$ large enough that
\[
 C_P\sqrt{3\cdot2^{13}C_\omega}
 (\ell_0+1)^{5/2}16^{-\ell_0}\leq1.
\]
Such an integer exists because the exponential factor dominates the
fixed power of $\ell_0+1$.  This proves
$\tau_\ell\sqrt{3V}\leq1$, with all the parameters recomputed from
this single choice of $\ell_0$.

Dyadic rounding of the rates makes no change to $V$, since $V$
depends on $\beta_i,A_i$ and $C_\omega$, not on the rates.
It preserves $\sum_i\eta_i^2\leq g^2$ and $\eta_{\max}\leq g$.
Consequently there is no factor-four enlargement of the curvature
budget from that rounding; its factor-two loss occurs only in the
final bound on $1/\eta_i$ in \eqref{eq:rates}.

Finally,
\[
 Y\leq\sqrt{3V}+\Lambda E\leq\sqrt{3V}+g\sqrt V E,
 \qquad\tau_\ell Y\leq1+gE.
\]
Together with \eqref{eq:prediction}, this gives
\[
 E\leq C_0+1+(C_1+1)gE\leq C_0+1+E/2,
 \qquad E\leq E_0:=2(C_0+1).
\]
The choice of $c$ implies $\eta_i^2E_0^2\leq1\leq A_i$.
Substituting into \eqref{eq:biasedenergy}, and using
$2\eta_i\rn{b_i}\leq A_i$, proves
$\eta_i\SwapReg_i(T)\leq4A_i$.  This is
\eqref{eq:inputconclusion} for the arbitrary chosen horizon.
\end{proof}

\subsection*{3. A rational normalizer with a small effective-score error}
\begin{lemma}\label{lem:normalization}
Given a rational row $z$, rational public parameters above, and rational
$0<\tau\leq1$, one can compute a positive rational probability row
$\widehat q$ satisfying
\[
 \widehat q=q_i(z+\zeta),\qquad\norm{\zeta}_\infty\leq\tau,
\]
for an analysis-only real vector $\zeta$.  Its bit complexity and output
bit length are polynomial in the input bit lengths, $s_i,d_i$, and
$\log(1/\tau)$.
\end{lemma}
\begin{proof}
Put $m=s_i$, $\beta=\beta_i$, $d=d_i$, and let $Z\geq1$ bound
$\norm z_\infty$.  The rational interval
\[
 \left[\min_bz_b-\frac3{m\beta},\ \max_bz_b\right]
\]
brackets the normalizing shift.  At the right endpoint all arguments
are nonpositive and
$\sum_bf_i(z_b-\mu)\leq mf_i(0)\leq em\beta/2
\leq e/2^{13}<1$, by the bound on $m\beta$ above.
At the left
endpoint all arguments are at least $3/(m\beta)$, where the defining
formula and Bernoulli's inequality give $f_i(v)>1/m$.
Take $R=2Z+3/(m\beta)+2$.
All relevant scalar arguments, and their unit neighborhoods, lie in
$[-R,R]$.

For $|v|\leq R$, $1/(R+1)\leq u(v),t_+(v)\leq R+1$.
The numerator in $f_i$ is at least $\beta$, and its denominator is
at most $3(R+1)^2$.  Hence
\[
 f_i(v)\geq\frac\beta{3(R+1)^2}.
\]
The logarithmic derivative of the numerator is
\[
 \frac{du}{(1+u^2)(1+du)}\geq\frac1{4(R+1)^4}.
\]
The negative logarithmic derivative of the denominator is nonnegative,
so
\begin{equation}\label{eq:derivativefloor}
 f_i'(v)\geq a_R:=\frac\beta{12(R+1)^6}.
\end{equation}
These deliberately conservative lower bounds are enough: their
logarithms grow only as $O(\log R+\log(1/\beta))$.

The function $f_i$ is composed of rational arithmetic, a square root,
and an integer power $d$.  On this interval every intermediate
magnitude and reciprocal of a positive denominator has logarithm
polynomial in $d,\log R,\log(1/\beta)$ and the input bit lengths.
For a direct certified implementation, bracket $\sqrt{v^2+4}$ by
rational bisection and evaluate the remaining expressions by rational
interval arithmetic.  One may use $t_+(v)=1/u(v)$ when this avoids a
subtraction of close values.  The positive denominators stay bounded
away from zero by the displayed estimates.  Repeated squaring or $d$
multiplications therefore evaluates $f_i$ to error $2^{-p}$ in a
number of bit operations polynomial in these quantities and $p$.

Certified bisection for the row sum locates its root to any prescribed
absolute error.  A nearly equal comparison does not require an
equality oracle: if an evaluation interval of width $2e$ contains one,
the derivative bound $|\sum_bf_i'(z_b-\mu)|\geq ma_R$ locates the
root within $2e/(ma_R)$ of that midpoint.  Otherwise the interval
certifies which half to retain.  The row sum also has an upper
derivative bound with polynomial logarithm.  Indeed $0<t_+'(v)\leq1$
and $|u'(v)|\leq1$; the logarithmic derivative of the numerator is
at most one and the negative logarithmic derivative of the denominator
is at most one.  Hence $f_i'(v)\leq2f_i(v)\leq2\beta(R+2)^d$,
and the row-sum derivative is at most $2m\beta(R+2)^d$.
This procedure computes approximations to the exact probability row
$q^*=q_i(z)$ with any prescribed absolute coordinate error, at the
claimed cost.

Compute positive dyadics $p_b$ with $|p_b-q_b^*|\leq e$ and put
$\widehat q_b=p_b/\sum_cp_c$.  If
\[
 e\leq\min\left\{\frac1{2m},\
 \frac{\beta}{12(R+1)^2},\
 \frac{a_R\tau}{8(m+1)}\right\},
\]
all coordinates are positive, the row is exactly normalized, and
\[
 |\widehat q_b-q_b^*|\leq2(m+1)e\leq a_R\tau/4.
\]
By \eqref{eq:derivativefloor}, $f_i$ changes by at least $a_R\tau$
on either length-$\tau$ side of the exact scalar argument.
Monotonicity therefore gives
\[
 \zeta_b=f_i^{-1}(\widehat q_b)-(z_b-\mu),\qquad
 |\zeta_b|\leq\tau.
\]
Now $\widehat q_b=f_i(z_b+\zeta_b-\mu)$ and its coordinates sum
to one.  Uniqueness of the normalizing shift proves the representation.
The inverse function and $\zeta$ are not computed.  All requested
accuracies have logarithm polynomial in the stated inputs, completing
the bit-complexity argument.
\end{proof}

\subsection*{4. Exact rational stationarity has polynomial bit cost}
\begin{lemma}\label{lem:stationarity}
A positive rational $m\times m$ stochastic matrix with at most $B$
bits in each numerator and denominator has an exactly normalized
rational stationary distribution computable with polynomially many
bit operations in $m,B$.  Its coordinate bit lengths are
$O(m^2B+m\log m)$.
\end{lemma}
\begin{proof}
Positivity implies existence and uniqueness of its stationary
distribution.  Replace one equation of $(Q^\top-I)y=0$ by
$\sum_by_b=1$; this gives a nonsingular rational system.  In each
equation, clearing denominators multiplies at most $m$ denominators,
so the resulting integer entries have $O(mB)$ bits.  The determinant
and every Cramer minor then have $O(m^2B+m\log m)$ bits, by the
determinant expansion or Hadamard's bound.  Fraction-free Gaussian
elimination, with row exchanges where necessary, uses polynomially
many arithmetic operations on integers of this size.  Cramer's rule
gives the output bound, and uniqueness gives exactly the required
stationary distribution.  This computation makes no numerical
condition-number assumption.
\end{proof}

This lemma is applied to a \emph{freshly formed} rational $Q$ in each
round.  Its denominators are never accumulated into future gain
sums; those gains are rounded to a common dyadic grid as follows.

\subsection*{5. The finite-bit learner}
Set $D=M+\ell$ and $B_*=10(D+1)^3$.  Choose an integer $b_0$ with
\begin{equation}\label{eq:b0}
 2^{-b_0}\leq\frac1{4096s(D+1)B_*^D},
 \qquad b_0=O(D\log(D+1)+\log s),
\end{equation}
and put
\[
 b_t=b_0+4\lceil\log_2(t+1)\rceil,\qquad h_t=2^{-b_t}.
\]
Then $h_t\leq 2^{-b_0}(t+1)^{-4}$ and $b_t$ is nondecreasing.

For each player maintain a cumulative gain matrix $C$, initially zero,
and the $D$ EMA state matrices $a_1,\ldots,a_D$ and three lag
matrices $c_1,c_2,c_3$, also initially zero.  Let
$\nu_j=1/(M+1)$ for $j\leq M$, and
$\nu_j=1/(\ell+1)$ for $j>M$.  At round $t$ perform:
\begin{enumerate}
\item Compute the rational predictor
\[
 \widetilde h=\sum_{j=1}^Da_j+3c_1-3c_2+c_3,
 \qquad\overline z=\eta_i(C+\widetilde h).
\]
\item For each row of $\overline z$, use Lemma~\ref{lem:normalization}
with $\tau=\eta_i/[4s_i(t+1)^2]$.  Form the positive rational
transition matrix $Q_i$ and compute its stationary distribution $y_i$
exactly using Lemma~\ref{lem:stationarity}.  Output the aggregated
original-action distribution $x_i$.
\item Request the original payoff vector at precision $b_t$ and copy
its coordinates to virtual actions.  Call the reply $\widetilde v_i$.
Set $\widetilde r_{a,b}$ to the nearest multiple of $2^{-b_t}$ to
$y_{i,a}(\widetilde v_{i,b}-\widetilde v_{i,a})$, using any fixed
tie convention.
\item Update $C\leftarrow C+\widetilde r$ \emph{exactly}, and set
$u\leftarrow\widetilde r$.  Sequentially for $j=1,\ldots,D$, set
\[
 u\leftarrow u-a_j,\qquad
 a_j\leftarrow\operatorname{round}_{b_t}(a_j+\nu_j u).
\]
The subtraction for $u$ is exact and uses the old $a_j$.
Finally update $(c_3,c_2,c_1)\leftarrow(c_2,c_1,u)$ simultaneously.
\end{enumerate}
All choices are deterministic.  Only public parameters, the player's
own state and payoff-vector replies are used.  Nothing requires $T$.

\subsection*{6. Rounding errors in the actual history}
Let $r_i^t$ be the \emph{true} gains of the strategies actually played
by this finite-bit learner, not the gains of an ideal comparison run.
Write
$d_i^t=\widetilde r_i^t-r_i^t$.  Oracle approximation and final gain
rounding give
\begin{equation}\label{eq:gainerror}
 \rn{d_i^t}\leq3s_i h_t.
\end{equation}
Indeed the payoff difference contributes at most $2y_{i,a}h_t$
in row $a$ and gain rounding contributes at most $h_t$ per row;
summing yields $(2+s_i)h_t\leq3s_i h_t$.
If $e_{i,j}^t$ denotes the rounding error in EMA update $j$, then
\begin{equation}\label{eq:emaerror}
 \rn{e_{i,j}^t}\leq s_i h_t.
\end{equation}

We now control the feedback of these errors into all later strategies.
For a scalar convolution kernel $K=\sum_{q\geq0}k_q\Sh^q$, define
\[
 \norm K_w=\sum_{q\geq0}(q+1)^2|k_q|,
 \qquad\norm w_{*,2}=\sup_{t\geq1}(t+1)^2\rn{w^t}.
\]
The inequality $t+1\leq(q+1)(t-q+1)$ gives
\begin{equation}\label{eq:weightedfilter}
 \norm{Kw}_{*,2}\leq\norm K_w\norm w_{*,2}.
\end{equation}
The same weight inequality makes these kernel norms submultiplicative.

For every stage, let $R_j=(I-(1-\nu_j)\Sh)^{-1}$ and
$F_j=R_j\Dif$.  Direct sums of the geometric kernel give
\[
 \norm{R_j}_w\leq2(D+1)^3,\qquad
 \norm{F_j}_w\leq B_*,\qquad
 \norm\Sh_w=4,\qquad\norm{\Dif^3}_w=56.
\]
For example $\sum_{q\geq0}(q+1)^2\rho^q
=(1+\rho)/(1-\rho)^3$, and $\norm{I-\Sh}_w=5$.
It follows that
\begin{equation}\label{eq:kernelbounds}
 \norm{I-\Pred}_w\leq57B_*^D,
 \qquad
 \norm{\Dif^3\prod_{h>j}F_h\Sh R_j}_w
 \leq56B_*^D.
\end{equation}

Let $U_j^t$ be the residual $u$ after stage $j$ in round $t$,
with $U_0=\widetilde r$.  The exact algebra of the rounded update is
\[
 U_j=F_jU_{j-1}-\Sh R_je_j.
\]
The stage telescoping identity and the three stored lags therefore
give the following \emph{exact} predictor identity:
\begin{align}
 \widetilde h_i&=(I-\Pred)r_i+\pi_i,\label{eq:predictorerror}\\
 \pi_i&=(I-\Pred)d_i+
 \Dif^3\sum_{j=1}^D\left(\prod_{h>j}F_h\right)\Sh R_je_{i,j}.
 \label{eq:piexact}
\end{align}
This calculation includes errors from every stage and every earlier
round.  In particular it does not identify the implemented trajectory
with an unrounded trajectory.

By \eqref{eq:gainerror}--\eqref{eq:kernelbounds},
\[
 \norm{\pi_i}_{*,2}
 \leq s_i(171+56D)B_*^D2^{-b_0}<\frac14.
\]
Also the infinite sum
\[
 b_i=\sum_{t=1}^\infty d_i^t
\]
converges absolutely, and \eqref{eq:b0} gives $\rn{b_i}\leq1/16$.
Its tail satisfies
\[
 \rn{\sum_{q=t}^\infty d_i^q}
 \leq s_i2^{-b_0}t^{-3}
 \leq\frac1{4(t+1)^2}.
\]
Here $\sum_{q=t}^\infty(q+1)^{-4}\leq1/(3t^3)$ and
$(t+1)^2/t^3\leq4$ for $t\geq1$.

The rational normalizer supplies an effective score perturbation
$\zeta_i^t$ with
$\rn{\zeta_i^t}/\eta_i\leq1/[4(t+1)^2]$.  Since
$C_i^{t-1}=\sum_{q<t}(r_i^q+d_i^q)$, its actual effective score is
\[
 z_i^t=\eta_i\left(b_i+\sum_{q<t}r_i^q
 +(I-\Pred)r_i^t+\varepsilon_i^t\right),
\]
where
\[
 \varepsilon_i^t=
 -\sum_{q=t}^\infty d_i^q+\pi_i^t+\frac{\zeta_i^t}{\eta_i},
 \qquad\rn{\varepsilon_i^t}\leq\frac3{4(t+1)^2}.
\]
Exact rational stationarity and the fixed virtual game give
$r_i^t=\Gamma_i(z^t)$.  These are precisely the hypotheses of
Lemma~\ref{lem:history}.

The infinite error sum $b_i$ is only a device of the proof.  The
learner does not compute it, store it, or use future replies.
Existence of the infinite run follows inductively from the finite,
well-defined algorithm at each round.  Deterministic pathwise
inequalities can then use the absolutely convergent error sum of
that run without any causal restriction on this analysis-only
constant.

\subsection*{7. Resource accounting and the regret conclusion}
\begin{theorem}\label{thm:transfer}
The learner above meets every resource requirement of prb_2a3012883e274a1384504ef3c9221892.
There are fixed polynomials $P,Q$
such that all its per-round bit costs, memory, requested precisions,
and output bit lengths are bounded by
$P(N,K,L,\lceil\log_2(t+1)\rceil)$, and
\[
 \sup_{T\geq1}\SwapReg_i(T)
 \leq\frac{8}{c}k^{5/2}\sqrt{A_iA}
 =O\!\left(\sqrt N\,K\log K\,\log^{5/2}(NK)\right)
 \leq Q(N,K,L).
\]
\end{theorem}
\begin{proof}
The preceding section verifies the hypotheses of Lemma~\ref{lem:history}
for the actual finite-bit run.  The lemma gives
$\SwapReg_i(T)\leq4A_i/\eta_i$;
\eqref{eq:rates} gives the first bound in the theorem, and
\eqref{eq:padding} transfers it to the original actions.
Since $s_i\leq2K$, $A_i\leq2K\log_2(2K)$,
$A\leq2NK\log_2(2K)$, and $k=O(\log(NK))$, the stated dimension
rate follows.  It is bounded by a fixed polynomial; for example a
sufficiently large universal multiple of $(NK)^6$ suffices.

For resources, $D=O(NK^2+\log(NK))$ and
$b_t=O(D\log(D+1)+\log(t+1))$.
The dyadic grid is nondecreasing.  Hence the exact cumulative gain
array $C$ has denominator dividing $2^{b_t}$, with integer part
$O(t)$: its bit length is $O(b_t+\log(t+1))$, rather than a sum of
all earlier precisions.  All EMA states and lag states are dyadics
on this grid.  A crude stable-filter bound using
$\norm{F_j}_{\mathrm{ker},1}\leq2$ bounds their magnitudes by
$2^{O(D)}$ times a polynomial in $D$; thus their integer parts require
only $O(D+\log(D+1))$ bits.  The rounding errors obey the same stable
linear recurrence, so the bound applies to the implemented states.
There are $O(Ds_i^2)$ such entries.

The raw score magnitude is at most $O(t+2^{O(D)})$, so the logarithm
of $R$ in Lemma~\ref{lem:normalization} is
$O(D+\log(t+1)+\log(NK))$.  Its tolerance obeys
\[
 \log(1/\tau)=O(\log(1/\eta_i)+\log s_i+\log(t+1)),
\]
which has the required polynomial bound.  Consequently each fresh
transition row has a polynomial bit bound, and the exact stationary
system has a polynomial bit cost and output size by
Lemma~\ref{lem:stationarity}.  Aggregating virtual coordinates can
use their common Cramer denominator, so it preserves a polynomial
output size and exact normalization.

The stationary distribution may have a different denominator in the
next round.  This causes no growth across time: products involving it
are rounded to the common dyadic grid before being added to $C$ or
passed to the EMA cascade.  Each temporary multiplication, division,
and rounding therefore has polynomial operand length.  The $
s_i^2D$ EMA operations and one own-payoff oracle query also have
polynomial bit cost and memory.  A single fixed polynomial can
dominate this finite composition of polynomial bounds, including
public initialization.  No part of the algorithm enumerates the
game's payoff table or stores past payoff vectors.
\end{proof}

\subsection*{8. Self-audit and verification status}
The proved numerical statements handle four separate issues:
positive rational row normalization at a small effective score;
exact rational stationarity with polynomial bit length;
bounded storage for cumulative gains; and the effect of all
feedback and EMA errors on the future score history.  The last issue
is handled by \eqref{eq:piexact} and the fixed-bias plus tail identity,
not by a bound on regret evaluation along an unchanged history.

The history extension was checked separately against the source's
first-time impulses, early and late time gates, score-window
comparison, and potential endpoints.  Lemma~\ref{lem:history}
records that extension as a proof using the cited analytic results,
with the energies, equation locators, and universal parameter bounds
specified.  The rate rounding preserves the general-rate hypotheses;
it does not silently use an identity restricted to the exact default
allocation.  All scalar inverse bounds are polynomial in the score
range, and the stationary-distribution denominators are prevented
from accumulating in future states by dyadic gain rounding.

The argument applies to every finite prefix of the same
horizon-independent run.  A one-action player is covered by virtual
duplicates and has zero original swap regret.  Positivity and exact
normalization hold before every stationary solve.  Regret throughout
the proof uses true expected gains of the implemented strategies.
These checks leave no identified gap in the new transfer proof.

The exact-real learning algorithm and its higher-order analytic
estimates remain credited to Mai.  This contribution uses those
estimates as cited mathematical results; it is not an independent
reconstruction of their entire proof.  Separate-agent review of a
proposed proof is not formal verification or external peer review.
No finite-precision learning implementation has been run, and no
numerical benchmark or experimental verification is claimed.

\section*{Models used}
OpenAI GPT-6 Astra Pro (system-reported model name; exact dated snapshot
unknown): source inspection, proof construction, and adversarial
self-review.  Separate agents checked the bounded-initial-score
extension and the numerical transfer, including resource accounting;
a coordinating agent reviewed the completed argument.  These were
model-based checks, not human or formal verification.

\section*{Human contributors and roles}
The user specified the problem and requested attempted solutions.
No human proof verification is claimed.  Tung Mai retains credit for
the exact-arithmetic algorithm and analytic estimates used here.
The question's originating research discussion credits Rida Laraki;
this statement does not attribute the present numerical construction
to him.

\begin{thebibliography}{9}
\raggedright
\bibitem{mai}
Tung Mai.
\emph{Constant Swap Regret in General-Sum Games via Two-Scale Higher-Order
Optimism}.  arXiv:2609.16751v2, revised 19 September 2026.
Sections 2--4; Appendix A.2--A.3; Lemmas 4.1--4.2, F.1 and G.1;
Appendix H.
\url{https://arxiv.org/html/2609.16751v2}.

\bibitem{statement}
We-Solve.
\emph{Finite-bit uncoupled learning with constant swap regret},
prb_2a3012883e274a1384504ef3c9221892, live submission key \texttt{ECON-3}.
Attached statement dated 17 October 2026; live research prompt
retrieved 18 October 2026.
\url{https://we-solve.github.io/problems/finite-bit-uncoupled-learning-with-constant-swap-regret/}.
\end{thebibliography}
\end{document}
